% !TEX root = main.tex
\section{양방향 추적과 전체 대상 동결}

\subsection{확정 시점과 감사 그래프}
감사 대상 그래프는 객체와 이벤트를 정점으로 갖는다. 실제 입
력 객체에서 소비 이벤트로, 그 이벤트에서 각 출력 객체로 방향
간선을 둔다. 각 입력의 생성 기록은 소비 기록보다 먼저 수락
되어야 한다. 한 번의 추적은 원장 시작부터 확정 시점까지 수
락된 불변 기록 구간(prefix) \(H_{t}\)를 기준으로 수행한다. \(H_{t}\)는 블
록체인·원장과 확정 블록 해시 또는 이에 대응하는 수락 순서를
특정하며, 미확정 블록 번호만을 확정성의 근거로 삼지 않는다.

모든 조회는 이 시점의 객체 존재, 생성 기록 색인, 기록의 이
벤트 형상, 출력 위치와 공개 트랜잭션 자료를 확인한다. 복호
화에는 외부 승인 절차를 거쳐 확보한 유효한 감사용 키와 기록
의 문맥을 사용한다. 자료 누락이나 발견된 디코딩·복호화·일관
성 오류는 실패로 처리한다. 부분 진단 자료를 남길 수 있지만
완성된 그래프나 전체 대상으로 간주해서는 안 된다.

\subsection{선행 이력 추적}
선행 추적 절차(TraceBackward)는 객체 또는 감사 기록 식별
자에서 시작한다. 객체 시작은 해당 객체와 생성 기록의 일치를
확인한 뒤 생성 이벤트로 이동한다. 기록 시작은 해당 이벤트
자체의 선행 경로를 조사하므로 출력이 없는 종료도 시작점이
될 수 있다.

각 방문 기록에서 암호화된 실제 부모 목록을 복호화하고, 부
모 수·타입·순서·중복 여부가 이벤트의 형상과 일치하는지 검사
한다. 각 부모의 생성 기록이 현재 기록보다 먼저 수락되었는지
확인하고, 부모에서 현재 이벤트로 가는 간선을 추가한다. 이후
각 부모의 생성 기록으로 이동하는 과정을 반복한다. 부모가 비
어 있으면 그 이벤트가 초기 등록인지 확인하고 시작 등록으로
기록한다.

같은 기록에 여러 경로로 도달하면 복호화와 확장을 한 번만
수행할 수 있지만, 새로 도달한 출력과 간선은 보존한다. 하나
의 출력에서 생성 기록을 읽었다는 이유로 그 생성 기록의 다른
출력까지 선행 경로에 추가하지 않는다. 결과는 복원한 선행 그
래프와 초기 등록 기록 집합이다. 이 알고리즘은 상태를 변경하
지 않는다.

\subsection{미래 소비 식별자를 이용한 후속 추적}
후속 추적 절차(TraceForward)가 객체에서 시작하면 그 객체
하나를 초기 큐에 넣고, 기록에서 시작하면 그 기록의 전체 출력
을 넣는다. 출력이 없는 기록을 지정하면 초기 객체 집합이 비
어 있다. 이를 종료 이전 경로의 조사로 해석하지 않는다.

도달 객체 \(r\)가 제품 주장이면 등록과 발급 이벤트의 생성 기
록을 확인하고 종단 제품 주장 집합에 넣는다. 회계 증표 또는
이전 증표이면 생성 기록에서 \(r\)의 출력 위치에 대응하는 미래
소비 암호문을 찾는다. 이를 복호화하여 타입별 소비 식별자
\(ν\)를 얻고, \(H_{t}\)에서 spentIn\([ν]\)를 조회한다. 그 값이 0이면 \(r\)는
추적이 멈추는 미소비 객체이며, 미소비 경계 집합(frontier)에
포함한다. 값이 aid이면 해당 이벤트가 소비 이벤트 후보이며,
생성 기록보다 나중의 기록인지와 공개 입력에 같은 타입의 \(ν\)가
포함되어 있는지를 확인한다.

소비 이벤트가 유효하면 \(r\)에서 그 이벤트로 가는 간선을 추
가한다. 처음 확장하는 소비 이벤트의 출력이 비어 있으면 종료

% Source PDF page 10, column 1
임을 확인하여 종료 집합에 넣는다. 그렇지 않으면 모든 출력을
그래프와 큐에 추가한다. 여러 입력 경로로 같은 소비 이벤트에
도달하더라도 각각의 입력 간선은 남기고 출력 확장은 한 번만
한다. 공개 소비 식별자와 실제 비공개 입력의 결합은 수락된
증명에 의존하므로, 이 단계에서 소비 이벤트의 다른 모든 부모
까지 복호화하여 역추적할 필요는 없다.

결과는 도달 그래프 \(G\), 회계 증표의 미소비 경계 집합 \(F_{N}\) ,
이전 증표의 미소비 경계 집합 \(F_{V}\) , 종단 제품 주장 집합 \(C\)와 종
료 기록 집합 \(X\)로 구분한다. \(F_{N},F_{V}\) 에는 객체의 타입 참조와
복호화한 소비 식별자를 함께 보관한다. 미소비 여부와 상태는
별개이므로 이미 동결 또는 사용 금지 상태인 객체도 미소비 경
계 집합에 포함한다.

\subsection{분기·합류에서의 도달 범위}
후속 범위는 시작 객체에서 소비 이벤트와 그 모든 출력을 따라
도달하는 디지털 그래프로 정의한다. 시작 객체가 병합에서 다
른 입력과 합쳐졌다면 병합 출력 전체를 추적한다. 분할과 가공
에서도 모든 출력을 포함하며, 폐기물 또는 질량 0이라는 이유
로 생략하지 않는다. 반면 병합의 다른 입력을 선행 방향으로
올라가서 같은 후속 집합에 넣지는 않는다.

시작 객체의 미래 식별자를 얻기 위해 생성 기록을 읽는 일과,
도달한 소비 이벤트에서 출력 전체로 확장하는 일도 구분한다.
전자만으로는 같은 생성 기록의 다른 출력을 추적 대상에 추가
하지 않는다. 따라서 그래프에 포함되는 범위는 단순히 읽은 기
록의 모든 객체가 아니라, 정해진 방향의 간선으로 실제 도달한
객체이다. 상충하는 객체·식별자 대응은 중복 도달로 숨기지 않
고 오류로 처리한다.

\subsection{차단 승인과 고정된 전체 대상}
감사·동결 절차(AuditAndFreeze)는 신뢰된 상태 관리자가 문
제 시작점의 차단을 승인한 뒤 실행하는 상위 절차이다. 오염
원인의 판정이나 선행 분석을 통한 시작점 선택은 선행 판단
이다. 모든 감사에 동결을 자동 적용하지 않으며, 조회 전용
TraceForward도 별도로 사용할 수 있다.

승인된 실행은 먼저 \(H_{t}\)에서 상태 변경 없이 후속 추적을 완
료한다. 추적 또는 미소비 경계 집합 일관성 검사가 실패하면 추
적 실패(TRACE\_FAILED)로 종료하고 이 실행의 동결 트랜잭션
은 한 건도 제출하지 않는다. 정상 결과를 얻으면 동결 대상을

\begin{equation}
    T=F_N\cup F_V,\quad t=(r,\nu)\tag{21}
\end{equation}
로 고정한다. 같은 타입 객체의 중복 도달은 하나로 정규화하
되, 동일 객체의 서로 다른 식별자 또는 서로 다른 객체의 동일
타입 식별자는 일관성 오류이다. 원시 식별자의 숫자가 같은 회
계 증표와 이전 증표의 식별자는 합치지 않는다.

\(T\) 는 미소비 경계 집합 전체와 정확히 같아야 한다. 이미 차
단된 객체, 폐기물, 질량 0 객체나 실행 중 실패한 객체도 임
의로 제외할 수 없다. 제품 주장과 종료는 감사 결과에 남지
만 동결 대상에는 넣지 않는다. \(T=\varnothing \)이면 미소비 대상 없음
(NO\_LIVE\_TARGETS)으로 종료한다. 이 결과는 해당 기준 시점
에서 동결할 회계 증표·이전 증표가 없다는 뜻이며, 문제의 부
재나 제품의 안전성을 판정하지 않는다.

% Source PDF page 10, column 2

\subsection{현재 상태를 이용한 동결 집행}
각 대상의 현재 소비 기록과 상태를 읽고 Table 4의 조치를 적
용한다. 미소비 활성 객체만 타입에 맞는 상태 변경으로 동결
상태에 옮긴다. 스마트 계약은 상태 관리자 권한, 식별자의 정
규형, 현재 미소비 여부와 허용 상태 전이를 검사한다. 그 식별
자가 특정 감사 암호문의 복호화 결과인지 또는 전체 미소비 경
계 집합이 포함되었는지를 스마트 계약이 별도의 증명으로 확
인하지는 않는다. 이는 신뢰된 감사자의 정확한 대상 선택과 집
행에 의존한다.

개별 동결 실패가 이전에 성공한 동결을 되돌리지는 않는다.
가능한 다른 대상의 처리는 계속하고, 공통 실행 장애로 시도하
지 못한 대상도 사유와 함께 결과에 남긴다. 시도·재시도·확인
은 유한한 한도 안에서 수행한다. 불명확한 제출 결과를 임의로
성공이나 실패로 판정하지 않는다.

트랜잭션 영수증과 최종 차단 상태는 구분한다. 예를 들어 다
른 승인된 실행이 먼저 동결하여 이번 제출이 거부되었더라도,
최종 조회에서 이미 차단된 상태로 확인될 수 있다. 반대로 과
거의 성공 영수증만으로 이후 상태까지 차단되었다고 보고해서
는 안 된다.

\subsection{공통 시점의 완료 판정}
집행 후 같은 원장의 확정 기록 구간 \(H_{r}\)을 선택하고, \(T\) 의 모든
객체를 그 시점에서 다시 조회한다. \(H_{r}\)은 \(H_{t}\) 이후 또는 동일한
기록 구간이며, 성공 근거인 트랜잭션과 실행에 사용한 상태 관
찰·처리 결과를 포함해야 한다. 서로 다른 시점에서의 개별 성
공을 모아 전체 동시 차단처럼 보고하지 않는다. 특히 소비 경
합을 관찰하고도 그 소비 이전의 기록 구간을 선택하여 완료로
보고할 수 없다.

대상 \(t\)의 타입을 \(κ_{t}\), 식별자를 \(ν_{t}\)라 하면

\begin{equation}
    \begin{aligned}\operatorname{Blocked}_{H_r}(t)\iff{}&\operatorname{spentIn}_{H_r}^{\kappa_t}[\nu_t]=0\\&\land\operatorname{status}_{H_r}^{\kappa_t}[\nu_t]\in\{\mathrm{Frozen},\mathrm{Revoked}\}.\end{aligned}\tag{22}
\end{equation}
공통 시점 완료(COMPLETE\_AT\_CHECKPOINT)는 다음 조건이
모두 참인 경우이다.

1. 후속 추적이 정상 완료되었고 \(T\) 가 비어 있지 않다.

2. \(H_{r}\)이 위 확정·시점 조건을 만족한다.

3. 보고된 대상 목록과 대상별 결과의 대상 집합이 각각 전체

\(T\) 와 같다.

4. 제출 여부 또는 처리 결과가 불명확한 트랜잭션이 없다.

5. 모든 \(t∈T\) 에 대해 식 (22)가 성립한다.

% Source PDF page 11, column 1
정상 추적 결과와 비어 있지 않은 \(T\) 를 얻었지만 이 조건을 만
족하지 못하면 미완료(INCOMPLETE)이다. 아직 활성인 대상,
소비 경합, 조회 실패, 미확정 제출과 미시도를 구분하여 기록한
다. 결과에는 시작점, \(H_{t}\), 추적 결과, 전체 \(T\) , 대상별 시도·영수
증·상태 관찰, 공통 \(H_{r}\) 또는 그 부재, 실행 전체의 최종 결과를
포함한다. 승인 부재나 잘못된 원장 문맥은 실행 전 거부한다.

\subsection{시점 사이의 소비와 후속 감사}
\(H_{t}\)에서 미소비 경계 집합에 속했던 객체가 동결 전에 소비되
면, 현재 소비 여부 검사 때문에 그 객체를 동결할 수 없다. 이를
\(T\) 에서 지우지 않고 미완료 사유로 남긴다. 이 상황은 과거 기준
시점의 추적 오류와 구분되지만, 해당 실행이 추가 사용 차단까
지 완료했다는 뜻도 아니다. 새 파생 객체를 차단하려면 새 기
준 시점에서 후속 감사를 수행해야 한다. 현재 \(T\) 를 새 집합으로
몰래 교체하거나 무한 재추적으로 완료를 보장하는 절차를 정
의하지 않는다.

또한 동결은 권한 있는 후속 조치로 해제될 수 있다. 완료 판
정은 \(H_{r}\)에서의 차단 상태에 관한 것이며, 그 상태가 유지되는
동안 후속 소비가 거부된다는 원장 규칙을 사용한다. 전체 절차
는 오프체인(off-chain) 감사와 여러 상태 트랜잭션의 연결이
다. 이를 하나의 원자적 온체인(on-chain) 트랜잭션 또는 새로
운 영지식증명 관계로 취급하지 않는다.


\begin{table}[htbp]\centering\small
    \setcounter{table}{3}\caption{Per-target action at execution time}
    \begin{tabularx}{\linewidth}{@{}XX@{}}\toprule
        \textbf{Observed state} & \textbf{Action in this run}          \\ \midrule
        Unspent, Active         & Submit a typed Freeze transaction    \\
        Unspent, Frozen/Revoked & Record an existing block             \\
        Already spent           & Retain target and consumer as a race \\
        Read failure            & Retain target as unconfirmed         \\
        \bottomrule\end{tabularx}\end{table}
